\documentclass[11pt,a4paper]{article}
\usepackage{kotex}
\usepackage{fontspec}
\setmainfont{Times New Roman}
\setmainhangulfont{AppleMyungjo}
\setsanshangulfont{Apple SD Gothic Neo}
\usepackage[a4paper,top=18mm,bottom=18mm,left=20mm,right=20mm]{geometry}
\usepackage{fancyhdr}
\setlength{\headheight}{24pt}
\usepackage{enumitem}
\usepackage{mdframed}
\usepackage{booktabs}
\linespread{1.18}
\setlength{\parindent}{0pt}
\pagestyle{fancy}
\fancyhf{}
\renewcommand{\headrulewidth}{0.6pt}
\fancyhead[L]{\sffamily\bfseries 2026 고1 영어 변형 - 정답 및 해설 지문 39}
\fancyhead[R]{\sffamily 삼림 파괴와 탄소 순환}
\fancyfoot[C]{\framebox[16mm][c]{\thepage}}
\newcommand{\qno}[1]{\par\vspace{10pt}{\sffamily\bfseries #1.}\ }
\newmdenv[linewidth=0.6pt,roundcorner=3pt,innertopmargin=7pt,innerbottommargin=7pt,
          innerleftmargin=9pt,innerrightmargin=9pt,skipabove=5pt,skipbelow=5pt]{pbox}
\begin{document}
\begin{center}
{\fontsize{15}{17}\selectfont\sffamily\bfseries [지문 39] 정답 및 해설}
\end{center}
\vspace{6pt}
\begin{center}
\renewcommand{\arraystretch}{1.35}
\begin{tabular}{cccc}
\toprule
문항 & 유형 & 정답 & 배점 \\
\midrule
1 & 문장 삽입 & ② & 3점 \\
2 & 요약문 & ③ & 2점 \\
3 & 영어 제목 & ① & 2점 \\
4 & 서술형 해석 & 모범답안 참조 & 2점 \\
\bottomrule
\end{tabular}
\end{center}

\qno{1}\textbf{문장 삽입 - 정답: ②}
\begin{pbox}
삽입 문장은 성숙한 나무 한 그루가 1년에 20kg 이상의 CO2를 흡수할 수 있다는 구체적 수치를 제시한다. 이는 바로 앞의 ``나무는 이산화탄소의 자연 조절자로 여겨진다''는 문장을 구체적으로 뒷받침한다. 따라서 ②가 가장 자연스럽다.

\medskip
①도 ``Trees play a huge role'' 뒤라서 매력적이지만, 그 뒤에는 광합성 설명이 이어져 나무의 역할을 먼저 원리 차원에서 설명한다. 수치 예시는 이 원리와 ``natural regulator''라는 평가가 나온 뒤에 들어가는 편이 더 매끄럽다. ③ 이후는 삼림 감소와 이산화탄소 증가로 논지가 넘어가므로 삽입문이 늦다.
\end{pbox}

\qno{2}\textbf{요약문 - 정답: ③}
\begin{pbox}
나무는 광합성을 통해 대기 중 CO2를 산소로 바꾸며, 더 많은 나무는 CO2 감소로 이어진다. 따라서 (A)는 \textit{removing}이 적절하다. 반대로 삼림 파괴는 이 자연적 과정을 막고 CO2 증가를 초래하므로 (B)는 \textit{disrupted}가 적절하다.

\medskip
①, ②의 \textit{producing}은 나무가 CO2를 만들어 낸다는 의미라 반대이다. ④ \textit{reinforced}와 ⑤ \textit{improved}는 삼림 파괴가 자연 과정을 강화하거나 개선한다는 뜻이 되어 본문과 반대이다.
\end{pbox}
\vspace{8pt}
\qno{3}\textbf{제목 - 정답: ①}
\begin{pbox}
글은 숲이 탄소 순환에서 CO2를 줄이는 자연 조절자 역할을 하며, 삼림 파괴가 이 과정을 방해해 지구 온난화를 심화시킨다고 설명한다. 따라서 ①이 가장 적절하다.
\end{pbox}

\qno{4}\textbf{서술형 해석}
\begin{pbox}
\textbf{모범답안}: 대기 중 이산화탄소가 더 많아지면, 태양 복사 에너지의 더 많은 부분이 우주로 방출되는 대신 지구로 다시 반사된다.

\medskip
\textbf{채점 기준}: ``more carbon dioxide in the atmosphere''를 대기 중 이산화탄소 증가로, ``sun's radiation''을 태양 복사 에너지로, ``reflected back to Earth instead of being released into space''를 우주로 방출되지 않고 지구로 다시 반사됨으로 옮기면 정답.
\end{pbox}

\end{document}
